\chapter{Implementierung}\label{chap:implementation}
    \section{Implementierung der Komponente MEDUSA}\label{sec:implementation-of-component-medusa}
        Die folgende Beschreibung der Implementierung von \ac{medusa} setzt die in \cref{subsec:component-medusa} entwickelte Konzeption in eine konkrete Realisierung in \texttt{C++} um.
        Soweit möglich wird auf die entsprechenden Abschnitte der Konzeption verwiesen, um redundante konzeptionelle Erläuterungen an dieser Stelle zu vermeiden.
        Der Fokus liegt stattdessen auf der konkreten Umsetzung, den getroffenen Implementierungsentscheidungen.
        \sig{Adrian Gabriel Axtmann}

        \subsection{Projektstruktur und Build-System}\label{subsec:impl-medusa-project-structure} \index{MEDUSA!Projektstruktur}
            Die Realisierung von \ac{medusa} erfolgt als \texttt{C++}-Projekt auf Basis von \gls{cmake} und wird gegen den Sprachstandard \texttt{C++23} kompiliert.
            Das Gesamtbild wird dabei von vier strukturellen Festlegungen geprägt.
            Damit spiegelt die Verzeichnisstruktur zunächst die in \cref{subsubsec:concept-medusa-modularization} konzipierte Schichtenarchitektur eins zu eins wider.
            Darauf abgestimmt organisiert das Build-System die Module als separate CMake-Targets mit explizit deklarierten Sichtbarkeiten.
            Zudem werden alle externen Bibliotheken deklarativ über \gls{fetchcontent} zur Konfigurationszeit gebunden.
            Plattformspezifische Anpassungen beschränken sich schließlich auf die reine Build-Konfiguration und die Initialisierung des \ac{OpenGL}-Kontexts.
            \sig{Adrian Gabriel Axtmann}

            \subsubsection{Verzeichnisstruktur}\label{subsubsec:impl-medusa-directory-layout}
                Die Verzeichnisstruktur des Repositories folgt der konzeptionellen Aufteilung in Infrastruktur-, Domänen-, Schnittstellen-, Präsentations- und Kompositionsschicht.
                Jedes Modul liegt in einem eigenen Unterverzeichnis unterhalb von \texttt{src/} und kapselt seinen öffentlichen Header sowie die zugehörige Implementierung.
                Tests, Beispieldaten, externe Abhängigkeiten und Shader sind aus dem Quellbaum ausgelagert, damit sich der reine Bibliothekscode in \texttt{src/} isoliert prüfen lässt.
                \sig{Adrian Gabriel Axtmann}

                \begin{table}[ht]
                    \centering
                    \caption[Top-Level-Verzeichnisse des Repositories]{Top-Level-Verzeichnisse des Repositories und ihre Zuordnung zur Schichtenarchitektur aus \cref{subsubsec:concept-medusa-modularization}}
                    \label{tab:impl-medusa-toplevel-dirs}
                    \renewcommand{\arraystretch}{1.2}
                    \begin{tabular}{@{} l l p{6.5cm} @{}}
                        \toprule
                        \textbf{Verzeichnis} & \textbf{Rolle} & \textbf{Inhalt} \\
                        \midrule
                        \texttt{src/}      & Bibliotheks- und Anwendungscode & Module der fünf Architekturschichten \\
                        \texttt{external/} & Abhängigkeiten                  & Per-Bibliothek-CMake-Fragmente und gevendortes GLAD \\
                        \texttt{tests/}    & Testsuite                       & Unit- und Integrationstests sowie Testdaten \\
                        \texttt{assets/}   & Laufzeitassets                  & GLSL-Shader für die Präsentationsschicht \\
                        \texttt{data/}     & Beispieldaten                   & Beispiel-\glspl{mesh} und Benutzereinstellungen \\
                        \texttt{ci/}       & CI-Infrastruktur                & Dockerfile \\
                        \bottomrule
                    \end{tabular}
                \end{table}
                \sig{Adrian Gabriel Axtmann}

            \subsubsection{Build-System-Konfiguration}\label{subsubsec:impl-medusa-cmake} \index{CMake}
                Als Build-System wird \gls{cmake} ab Version~3.22.1 eingesetzt.
                Die oberste Ebene legt die Sprachversion auf \texttt{C++23} fest, deaktiviert nicht-portable Compiler-Erweiterungen und bindet die in \cref{tab:impl-medusa-cmake-options} aufgeführten Optionen ein.

                \begin{table}[ht]
                    \centering
                    \caption[CMake-Optionen des Wurzelprojekts]{CMake-Optionen des Wurzelprojekts und ihre Voreinstellungen}
                    \label{tab:impl-medusa-cmake-options}
                    \renewcommand{\arraystretch}{1.2}
                    \begin{tabular}{@{} l l p{6.0cm} @{}}
                        \toprule
                        \textbf{Option} & \textbf{Standard} & \textbf{Wirkung} \\
                        \midrule
                        \texttt{MEDUSA\_BUILD\_APP}   & \texttt{ON}    & Baut die grafische Anwendung \\
                        \texttt{MEDUSA\_BUILD\_TESTS} & \texttt{ON}    & Baut die Testsuite \\
                        \texttt{CMAKE\_BUILD\_TYPE}   & \texttt{Debug} & Build-Typ \\
                        \texttt{FETCH\_USE\_SHALLOW}  & \texttt{ON}    & Flache Git-Clones für FetchContent \\
                        \bottomrule
                    \end{tabular}
                \end{table}

                Jedes Modul der Schichtenarchitektur wird als eigenständiges statisches Bibliotheks-Target angelegt und über einen Alias mit \texttt{medusa::}-Präfix exportiert.
                Die Linksichtbarkeit ist konsequent so gewählt, dass nur die für die öffentliche API benötigten Abhängigkeiten als \texttt{PUBLIC} und alle übrigen als \texttt{PRIVATE} propagiert werden.
                \sig{Adrian Gabriel Axtmann}

                \begin{lstlisting}[style=cmakeStyle, caption={Linkkonfiguration von \texttt{medusa::slicing} (Ausschnitt).}, label={lst:impl-medusa-slicing-cmake}]
add_library(medusa_slicing STATIC ${SLICING_SOURCES})
add_library(medusa::slicing ALIAS medusa_slicing)

target_link_libraries(medusa_slicing PUBLIC
    medusa::core
    medusa::geometry)

# Crystal-spezifische Abhaengigkeiten bleiben PRIVATE
target_link_libraries(medusa_slicing PRIVATE
    ext::eigen
    ext::libigl)
                \end{lstlisting}

                Externe Bibliotheken werden über \gls{fetchcontent} aus ihren öffentlichen Git"=Repositories bezogen.
                Pro Bibliothek existiert eine eigene CMake-Datei, die das Versions-Tag, Repository-\ac{URL} sowie Konfigurationsoptionen kapselt.
                Eine einzige Ausnahme bildet GLAD, das als gevendorter Quellbaum mitgeführt wird, da der offizielle Generator eine Webanfrage zur Build-Zeit stellen würde.
                \sig{Adrian Gabriel Axtmann}

            \subsubsection{Plattformspezifische Anpassungen}\label{subsubsec:impl-medusa-platform-specifics}
                Die in \cref{subsec:architecture-principles-and-design-goals} geforderte Plattformunabhängigkeit wird durch plattformneutrale Bibliotheken (GLFW, GLM, Assimp, Eigen) abgedeckt.
                Zwei Stellen bleiben dennoch plattformsensitiv.

                \paragraph{Compiler-Anforderungen}\label{para:impl-medusa-platform-compilers}
                    Die Mindestversionen sind GCC~13, Clang~16 sowie MSVC~19.35 (Visual Studio~2022~17.5).
                    Diese Schwellen werden in der \texttt{CMakeLists.txt} geprüft und mit einer Fehlermeldung quittiert.
                    \sig{Adrian Gabriel Axtmann}

                \paragraph{OpenGL-Kontext auf macOS}\label{para:impl-medusa-platform-macos-gl}
                    Apple stellt seit macOS~10.14 OpenGL ausschließlich im Forward-Compatible Core-Profile bereit und unterstützt maximal Version~4.1.
                    Der OpenGL-Kontext wird daher durchgängig als \texttt{3.2 Core, Forward-Compatible} angefordert und alle Shader werden in GLSL~150 geschrieben.
                    \sig{Adrian Gabriel Axtmann}

        \subsection{Infrastruktur-Schicht}\label{subsec:impl-medusa-infrastructure}
            Die Infrastruktur-Schicht stellt querschnittliche Funktionalität bereit, auf die alle anderen Schichten zugreifen.
            Das Modul \texttt{core} enthält die Klasse \texttt{Logger}, die als dünner Wrapper um \texttt{spdlog} fungiert.
            \sig{Adrian Gabriel Axtmann}

            \subsubsection{Logging-Subsystem}\label{subsubsec:impl-medusa-logging}\index{Logging}\index{spdlog}
                Drei Designentscheidungen prägen das Subsystem.
                So ist der Logger grundsätzlich global zugänglich.
                Zudem arbeitet er asynchron, sodass schreibende Threads nicht durch I/O blockiert werden.
                Das Absetzen der Logmeldungen erfolgt schließlich über Makros.
                \sig{Adrian Gabriel Axtmann}

                \paragraph{Sink-Konfiguration}\label{para:impl-medusa-logger-sinks}
                    Es werden parallel zwei Sinks bedient: ein farbig formatierter Konsolen-Sink auf \texttt{stderr} und ein rotierender Dateisink.
                    Der Konsolen-Sink filtert standardmäßig auf \texttt{info}, der Datei-Sink auf \texttt{trace}.
                    \sig{Adrian Gabriel Axtmann}

                \paragraph{Compile-Time-Filterung}\label{para:impl-medusa-logger-compile-time}
                    Über die Präprozessor-Konstante \texttt{SPDLOG\_ACTIVE\_LEVEL} werden Logaufrufe unterhalb der angegebenen Stufe bereits zur Übersetzungszeit aus dem Maschinencode entfernt.
                    In Release-Übersetzungen werden \texttt{trace}- und \texttt{debug}-Aufrufe eliminiert.
                    \sig{Adrian Gabriel Axtmann}

                    \begin{lstlisting}[style=cppStyle, caption={Verwendung der Logger-Makros.}, label={lst:impl-medusa-logger-macros}]
Logger::init("logs");
MEDUSA_INFO("Startup complete");
MEDUSA_DEBUG("Loaded {} vertices, {} faces", n_vert, n_face);
MEDUSA_WARN("Mesh contains {} non-manifold edges", n_bad);
Logger::shutdown();
                    \end{lstlisting}

        \subsection{Domänen-Schicht: Geometrieverarbeitung}\label{subsec:impl-medusa-geometry}\index{MEDUSA!Modul!geometry}
            Die Domänen-Schicht implementiert die in \cref{subsubsec:concept-medusa-data-representations} beschriebene indizierte Dreiecksrepräsentation.
            Sie ist bewusst so geschnitten, dass sie keine Datei-IO-Abhängigkeit besitzt und keinerlei Bezug zur grafischen Darstellung kennt.
            \sig{Adrian Gabriel Axtmann}

            \subsubsection{Datenstrukturen für Mesh}\label{subsubsec:impl-medusa-mesh-data-structures}\index{Mesh!Datenstruktur}
                Die zentrale Datenstruktur \texttt{TriangleMesh} verwendet ein \ac{SoA}-Layout, bei dem \glspl{vertex}, Normalen, Flächenindizes sowie Adjazenz- und Inzidenzlisten in separaten \texttt{std::vector}-Instanzen liegen.
                Diese Wahl verbessert die Cache-Auslastung bei Iterationen über Teilmengen der Attribute.

                \begin{lstlisting}[style=cppStyle, caption={Indizierte Dreiecksrepräsentation \texttt{TriangleMesh}.}, label={lst:impl-medusa-trianglemesh}]
struct TriangleMesh
{
    std::vector<glm::vec3>             vertices;
    std::vector<glm::vec3>             normals;
    std::vector<glm::uvec3>            faces;
    std::vector<glm::vec3>             face_normals;
    std::vector<glm::uvec3>            face_to_face;
    std::vector<std::vector<uint32_t>> vertex_to_faces;
    AABB                               bounds;

    static constexpr uint32_t NO_ADJACENT_FACE
        = std::numeric_limits<uint32_t>::max();

    bool isValid() const noexcept;
};
                \end{lstlisting}

                Die Flächen-zu-Flächen-Adjazenz ist als \texttt{glm::uvec3} pro \gls{triangle} kodiert, wobei Boundary-\glspl{edge} mit \texttt{NO\_ADJACENT\_FACE} markiert werden.
                Die Inzidenz ist als \texttt{std::vector<std::vector<uint32\_t>>} repräsentiert, da die Anzahl inzidenter \glspl{triangle} pro \gls{vertex} nicht beschränkt ist.
                \sig{Adrian Gabriel Axtmann}

            \subsubsection{Mesh-Import über Assimp}\label{subsubsec:impl-medusa-mesh-import}
                \index{Mesh!Import}
                Der Mesh-Import erfolgt über die Klasse \texttt{MeshLoader}, die einen \texttt{Assimp::Importer} mit fester Postprocess-Konfiguration kapselt.

                \begin{table}[ht]
                    \centering
                    \caption[Verwendete Assimp-Postprocess-Flags]{%
                        Verwendete Postprocess-Flags des Assimp-Importers%
                    }
                    \label{tab:impl-medusa-assimp-flags}
                    \renewcommand{\arraystretch}{1.2}
                    \begin{tabular}{@{} l p{7.5cm} @{}}
                        \toprule
                        \textbf{Flag} & \textbf{Wirkung} \\
                        \midrule
                        \texttt{aiProcess\_Triangulate}        & Zerlegt nicht-dreieckige Polygone \\
                        \texttt{aiProcess\_GenSmoothNormals}   & Erzeugt geglättete Vertex-Normalen \\
                        \texttt{aiProcess\_JoinIdenticalVertices} & Verschmilzt identische Vertizes \\
                        \texttt{aiProcess\_ImproveCacheLocality}  & Reordert Faces für Cache-Auslastung \\
                        \texttt{aiProcess\_PreTransformVertices} & Wendet Szenen-Hierarchie an \\
                        \bottomrule
                    \end{tabular}
                \end{table}

                Mehrere Submeshes werden vor der Konvertierung zusammengeführt.
                Material- und Texturkoordinaten werden hierbei verworfen.
                \sig{Adrian Gabriel Axtmann}

            \subsubsection{Vertex-Verschmelzung und Adjazenzaufbau}\label{subsubsec:impl-medusa-vertex-welding}
                \index{Vertex Welding}
                Das positionsbasierte Verschmelzen der Vertizes ist Voraussetzung dafür, dass nachgelagerte Stufen über die Adjazenzlisten navigieren können.
                Die Verschmelzung erfolgt über eine \texttt{std::unordered\_map} mit quantisiertem \texttt{glm::vec3} als Schlüssel.
                Die Quantisierung rundet die Komponenten auf ein Gitter mit Maschenweite $10^{-5}\,\si{mm}$.

                \begin{lstlisting}[style=cppStyle, caption={Quantisierter Hash für Vertex-Welding.}, label={lst:impl-medusa-vertex-hash}]
struct QuantizedVec3Hash
{
    static constexpr float weld_eps = 1.0e-5f;

    std::size_t operator()(const glm::vec3& v) const noexcept
    {
        const auto qx = static_cast<int32_t>(std::round(v.x / weld_eps));
        const auto qy = static_cast<int32_t>(std::round(v.y / weld_eps));
        const auto qz = static_cast<int32_t>(std::round(v.z / weld_eps));

        std::size_t h = std::hash<int32_t>{}(qx);
        h ^= std::hash<int32_t>{}(qy) + 0x9e3779b9u + (h << 6) + (h >> 2);
        h ^= std::hash<int32_t>{}(qz) + 0x9e3779b9u + (h << 6) + (h >> 2);
        return h;
    }
};
                \end{lstlisting}

                Nach der Verschmelzung baut \texttt{buildVertexToFaces()} die Inzidenzlisten auf, anschließend bestimmt \texttt{buildFaceAdjacency()} die Flächen-zu-Flächen-Adjazenz.
                Der Adjazenzaufbau ist erwartet linear in der Anzahl der \glspl{triangle}.
                \sig{Adrian Gabriel Axtmann}

            \subsubsection{Normalen und Bounding-Box}\label{subsubsec:impl-medusa-normals-aabb}
                \index{Bounding Box}
                Face-Normalen werden pro \gls{triangle} als normiertes Kreuzprodukt berechnet.
                Vertex-Normalen werden vom Assimp-Importer mit \texttt{aiProcess\_GenSmoothNormals} erzeugt.
                Die \ac{AABB} wird in einem Pass über das \texttt{vertices}-Array aufgebaut.
                Sie wird von den Slicing-Algorithmen zur Schichtanzahl-Berechnung und von der Visualisierung zur Kameraplatzierung verwendet.
                \sig{Adrian Gabriel Axtmann}

        \subsection{Schnittstellen-Schicht: Toolpath-Export}\label{subsec:impl-medusa-exporter} \index{MEDUSA!Modul!exporter}
            Die Schnittstellen-Schicht setzt die in \cref{subsubsec:interface-conceptual-structure} spezifizierte Toolpath"=Repräsentation in ein konkretes \gls{json}-Schema um.
            \sig{Adrian Gabriel Axtmann}

            \subsubsection{JSON-Schema und Serialisierung}\label{subsubsec:impl-medusa-json-schema} \index{JSON-Schema}
                Die POD-Strukturen \texttt{JobMetadata}, \texttt{Waypoint} und \texttt{Toolpath} bilden Header und Wegpunktstrom ab.
                Die Position ist als drei \texttt{double}-Komponenten in \si{mm} repräsentiert, die Orientierung als \gls{quaternion} mit Reihenfolge $(q_w, q_x, q_y, q_z)$.
                Intern speichert die \texttt{Segment}-Struktur die Ausrichtung lediglich als \texttt{glm::vec3} (entspricht $\vec{z}_\mathrm{tool}$).
                Die Konvertierung in ein vollständiges \gls{quaternion} ist noch nicht implementiert.


                \begin{lstlisting}[style=cppStyle, caption={POD-Struktur \texttt{Waypoint}.}, label={lst:impl-medusa-waypoint}]
struct Waypoint
{
    double x, y, z;                 // Position in mm
    double qw, qx, qy, qz;          // Einheits-Quaternion
    double feed_rate;               // mm/s
    double extrusion;               // mm^3/s, 0 fuer Travel-Moves
    MotionType    motion_type;      // Travel | Print
    std::uint32_t layer_id;
    std::uint32_t segment_id;
    std::uint32_t branch_id;
};
                \end{lstlisting}

                Die Serialisierung erfolgt über \texttt{nlohmann\_json} mit \ac{ADL}-basierten \texttt{to\_json}\slash\allowbreak\texttt{from\_json}"=Funktionen.
                Aufzählungstypen werden als Kleinbuchstaben-Zeichenkette serialisiert.
                \sig{Adrian Gabriel Axtmann}

            \subsubsection{Prüfsumme}\label{subsubsec:impl-medusa-checksum}\index{FNV-1a-Hash}
                Die Prüfsumme über den Wegpunktstrom verwendet \gls{FNV-1a}-64, einen nicht-kryptografischen 64"=Bit"=Hash.
                Der Hash wird über die serialisierte \gls{json}"=Repräsentation gebildet, nicht über die Rohdaten.
                Dadurch kann \ac{kronos} die Prüfsumme unabhängig nachrechnen: es genügt, das empfangene \texttt{waypoints}-Array erneut zu serialisieren und zu hashen, ohne die interne \texttt{Waypoint}-Struktur zu kennen.
                \sig{Adrian Gabriel Axtmann}

                \begin{lstlisting}[style=cppStyle, caption={FNV-1a-64-Implementierung.}, label={lst:impl-medusa-fnv1a}]
constexpr std::uint64_t k_fnv1a_offset = 14695981039346656037ULL;
constexpr std::uint64_t k_fnv1a_prime  = 1099511628211ULL;

std::uint64_t fnv1a_64(std::string_view data) noexcept
{
    std::uint64_t hash = k_fnv1a_offset;
    for (unsigned char c : data)
    {
        hash ^= static_cast<std::uint64_t>(c);
        hash *= k_fnv1a_prime;
    }
    return hash;
}
                \end{lstlisting}

            \subsubsection{Validierung und atomares Schreiben}\label{subsubsec:impl-medusa-atomic-write} \index{Validator}
                Vor jedem Schreibvorgang prüft die Klasse \texttt{Validator} fünf Eigenschaften: Metadaten-Vollständigkeit, nichtleerer Wegpunktstrom, finite Werte, \gls{quaternion}-Norm ($\lVert q \rVert^2 \approx 1$ mit Toleranz $10^{-6}$) und monotone Layer-IDs.
                Eine kinematische Erreichbarkeitsprüfung ist über das Interface \texttt{IReachabilityCheck} pluggable.

                Das Schreiben erfolgt zweistufig: vollständige \gls{json}"=Repräsentation in temporäre Datei mit Suffix \texttt{.tmp}, dann atomares Umbenennen über \texttt{std::filesystem::rename()}.
                Der Dateiname folgt dem Muster \texttt{toolpath\_<UTC-Timestamp>\_<algorithm>.json}.

                Die Facade \texttt{KronosExporter} verwendet einen eigenständigen \texttt{Result<T, E>}-Typ auf Basis von \texttt{std::variant}, da das Modul als eigenständige Bibliothek auch in Projekten ohne \texttt{C++23} verwendbar bleiben soll.
                \sig{Adrian Gabriel Axtmann}

        \subsection{Präsentations-Schicht: 3D-Visualisierung}\label{subsec:impl-medusa-graphics}\index{MEDUSA!Modul!graphics}
            Die Visualisierung nutzt \ac{OpenGL}~3.2 Core Profile und ist nach dem Prinzip \glqq eine Renderschicht pro Klasse\grqq{} organisiert.
            Dabei ist jede visuelle Darstellungsebene --- Mesh, Toolpath, Hilfsgeometrie, Feld-Overlay --- einer eigenen Renderer-Klasse zugeordnet, die ausschließlich die \ac{VAO}/\ac{VBO}-Ressourcen und das Shader-Paar ihrer Schicht verwaltet.
            Schichten werden unabhängig aktiviert und können so zur Laufzeit ein- und ausgeblendet werden.
            \sig{Adrian Gabriel Axtmann}

            \subsubsection{OpenGL-Kontext}\label{subsubsec:impl-medusa-opengl-context}\index{OpenGL!Kontextinitialisierung}
                Der \ac{OpenGL}-Kontext wird zweistufig erzeugt: GLFW legt das Fenster mit Kontext-Hints an, GLAD lädt die Funktionszeiger.
                \sig{Adrian Gabriel Axtmann}

                \begin{lstlisting}[style=cppStyle, caption={GLFW-Kontext-Hints für einheitliches OpenGL-Profil.}, label={lst:impl-medusa-glfw-hints}]
glfwWindowHint(GLFW_CONTEXT_VERSION_MAJOR, 3);
glfwWindowHint(GLFW_CONTEXT_VERSION_MINOR, 2);
glfwWindowHint(GLFW_OPENGL_PROFILE,        GLFW_OPENGL_CORE_PROFILE);
glfwWindowHint(GLFW_OPENGL_FORWARD_COMPAT, GL_TRUE);   // macOS
                \end{lstlisting}

            \subsubsection{Renderer-Klassen}\label{subsubsec:impl-medusa-renderer-classes}\index{Renderer-Klassen}
                \begin{table}[ht]
                    \centering
                    \caption[Renderer-Klassen des Graphics-Moduls]{%
                        Renderer-Klassen und ihre Zuordnung zu Renderschicht und Shader-Paar
                    }
                    \label{tab:impl-medusa-renderers}
                    \renewcommand{\arraystretch}{1.2}
                    \begin{tabular}{@{} l l l @{}}
                        \toprule
                        \textbf{Klasse} & \textbf{Renderschicht} & \textbf{Shader-Paar} \\
                        \midrule
                        \texttt{SceneRenderer}        & Mesh-Schicht           & \texttt{scene.vert/.frag} \\
                        \texttt{AxesRenderer}         & Hilfsgeometrie         & \texttt{axes.vert/.frag} \\
                        \texttt{GridRenderer}         & Hilfsgeometrie         & \texttt{grid.vert/.frag} \\
                        \texttt{ToolpathRenderer}     & Toolpath-Schicht       & \texttt{toolpath.vert/.frag} \\
                        \texttt{PhiOverlayRenderer}   & Feld-Overlay           & \texttt{phi\_overlay.vert/.frag} \\
                        \bottomrule
                    \end{tabular}
                \end{table}

                Auf eine formale Basisklasse wurde verzichtet.
                Jede Renderer-Klasse verpflichtet sich auf drei Methoden: Konstruktor (Shader anlegen), \texttt{render(...)} und Destruktor (\ac{GPU}-Ressourcen freigeben).
                Alle Shader verwenden GLSL~150 und feste Attributpositionen (\texttt{aPos} auf Location~0, \texttt{aNormal} auf Location~1).
                \sig{Adrian Gabriel Axtmann}

            \subsubsection{Viridis-Farbabbildung}\label{subsubsec:impl-medusa-viridis-shader}\index{Viridis}
                Das $\Phi$-Overlay verwendet eine polynomielle Approximation der Viridis-Palette direkt im Fragment-Shader.
                \sig{Adrian Gabriel Axtmann}

                \begin{lstlisting}[style=glslStyle, caption={Viridis-Farbabbildung als polynomielle Approximation.}, label={lst:impl-medusa-viridis-frag}]
vec3 viridis(float t)
{
    const vec3 c0 = vec3( 0.2777, 0.0054, 0.3340);
    const vec3 c1 = vec3( 0.1056, 1.4046, 1.3845);
    const vec3 c2 = vec3(-0.3308, 0.2148, 0.0952);
    const vec3 c3 = vec3(-4.6342,-5.7991,-19.3324);
    const vec3 c4 = vec3( 6.2289, 14.1799, 56.6905);
    const vec3 c5 = vec3( 4.7763, -13.7451, -65.3530);
    const vec3 c6 = vec3(-5.4354, 4.6458, 26.3124);
    return c0 + t*(c1 + t*(c2 + t*(c3 + t*(c4 + t*(c5 + t*c6)))));
}
                \end{lstlisting}


            \subsubsection{Daten-Upload-Strategie}\label{subsubsec:impl-medusa-buffer-updates}\index{VBO!Upload}
                Bei Änderungen wird der Puffer vollständig über \texttt{glBufferData()} neu befüllt.
                Eine \ac{GPU}-Aktualisierung erfolgt in drei Situationen: beim Laden eines neuen Meshes, nach Abschluss eines Slicing-Laufs und beim Umschalten zwischen Renderschichten.
                Sämtliche \ac{GPU}-Ressourcen werden über \ac{RAII} verwaltet.
                \sig{Adrian Gabriel Axtmann}

        \subsection{Präsentations-Schicht: Bedienoberfläche}\label{subsec:impl-medusa-ui}\index{MEDUSA!Modul!ui}
            Die Bedienoberfläche basiert auf Dear ImGui.
            Das Docking-Feature aus dem \texttt{docking}-Branch ermöglicht bewegliche Werkzeugfenster.
            Position, Größe und Dock"=Zuweisung werden über \texttt{imgui.ini} zwischen Sitzungen persistiert.
            \sig{Adrian Gabriel Axtmann}

            \subsubsection{Werkzeugfenster}\label{subsubsec:impl-medusa-tool-windows}\index{ImGui!Werkzeugfenster}
                Drei Werkzeugfenster realisieren die Interaktionsmuster.
                \texttt{MeshFileBrowser} scannt ein konfiguriertes Verzeichnis nach 3D-Dateien und stellt sie als auswählbare Liste dar; ein Klick lädt das gewählte Modell.
                Das Parameter-Panel zeigt Eingabefelder für den Slicer wie bspw. Schichtdicke.
                Das Diagnose-Panel zeigt nach jedem Slicing-Lauf Schichtanzahl, Segmentanzahl, Branch-Anzahl, Coverage und Laufzeit.

                Die Aktionsschaltflächen \emph{Slice}, \emph{Visualize} und \emph{Export} in der Menüleiste lösen die Pipeline-Stufen explizit aus.
                \sig{Adrian Gabriel Axtmann}

        \subsection{Präsentations-Schicht: Eingabe und Kamera}\label{subsec:impl-medusa-controls}\index{MEDUSA!Modul!controls}
            Das Modul \texttt{controls} übersetzt Maus- und Scrollrad-Eingaben in Zustandsänderungen.
            \texttt{core::Camera} ist als Orbit-Kamera modelliert mit Azimut, Elevation und Distanz.
            \sig{Adrian Gabriel Axtmann}

            \begin{table}[ht]
                \centering
                \caption[Eingabe-Aktion-Zuordnung]{Zuordnung der Eingaben zu Aktionen}
                \label{tab:impl-medusa-input-mapping}
                \renewcommand{\arraystretch}{1.2}
                \begin{tabular}{@{} l l l @{}}
                    \toprule
                    \textbf{Eingabe}             & \textbf{Aktion}        & \textbf{Ziel} \\
                    \midrule
                    Rechte Maustaste + Drag      & Azimut und Elevation   & \texttt{Camera} \\
                    Scrollrad                    & Distanz                & \texttt{Camera} \\
                    Linke Maustaste + Drag       & Yaw und Pitch          & \texttt{ModelTransform} \\
                    \bottomrule
                \end{tabular}
            \end{table}

            Die Trennung von \ac{UI}- und Viewport"=Eingaben erfolgt über \texttt{ImGuiIO::\allowbreak WantCaptureMouse}.
            Beide Controller brechen ihre Verarbeitung ab, wenn die \ac{UI} die Eingabe beansprucht.
            \sig{Adrian Gabriel Axtmann}

        \subsection{Komposition: Anwendungsschicht}\label{subsec:impl-medusa-app}\index{MEDUSA!Modul!app}
            Die Anwendungsschicht orchestriert alle Module und ist die einzige Schicht mit Kenntnis aller anderen.
            Das Modul besteht aus \texttt{main.cpp} (Prozess-Einsprungspunkt) und \texttt{app.cpp} (Orchestrator-Klasse).
            \sig{Adrian Gabriel Axtmann}

            \subsubsection{Lebenszyklus und Hauptschleife}\label{subsubsec:impl-medusa-main-loop}
                Der Lebenszyklus umfasst drei Phasen.
                Die Initialisierungsphase erzeugt \ac{OpenGL}-Kontext, ImGui, Renderer und Controller in fester Reihenfolge.
                Die Hauptschleife folgt dem klassischen Aufbau: Eingaben einsammeln, Controller aktualisieren, 3D-Szene rendern, ImGui-Overlay zeichnen, Puffer tauschen.
                Das Frame-Pacing folgt V-Sync über \texttt{glfwSwapInterval(1)}.
                Die Beendigungsphase läuft in umgekehrter Reihenfolge.

                \begin{lstlisting}[style=cppStyle, caption={Aufbau der Hauptschleife.}, label={lst:impl-medusa-loop}]
while (!glfwWindowShouldClose(window_))
{
    glfwPollEvents();

    const bool ui_capture = ImGui::GetIO().WantCaptureMouse;
    camera_ctrl_.updateFromMouseDrag(window_, ui_capture);
    model_ctrl_.updateFromMouseDrag(window_,  ui_capture);

    glViewport(0, 0, fb_w_, fb_h_);
    glClear(GL_COLOR_BUFFER_BIT | GL_DEPTH_BUFFER_BIT);

    renderScene3d();
    imgui_.beginFrame();
    renderImGui();
    imgui_.endFrameAndRender();

    glfwSwapBuffers(window_);
}
                \end{lstlisting}
                \sig{Adrian Gabriel Axtmann}

            \subsubsection{Pipeline-Verkettung}\label{subsubsec:impl-medusa-pipeline-wiring}
                \index{App!Pipeline-Verkettung}
                Die Pipeline-Verkettung erfolgt in drei Methoden.
                \texttt{loadMesh()} lädt das Mesh über \texttt{MeshLoader}, konvertiert in \texttt{TriangleMesh} und richtet die Kamera auf die \ac{AABB} aus.
                \texttt{runSlicing()} delegiert an \texttt{SlicingPipeline::execute()} und übergibt das Ergebnis an die Renderer.
                \texttt{exportToolpath()} reicht den Toolpath an \texttt{KronosExporter::\allowbreak export\_job()} weiter.
                Die Visualisierungs-Abzweigung ist read-only durch \texttt{const}-Referenzen realisiert.
                \sig{Adrian Gabriel Axtmann}

        \subsection{Test- und CI-Setup}\label{subsec:impl-medusa-testing}
            \index{MEDUSA!Testkonzept}
            Die Tests verwenden GoogleTest und sind in \texttt{tests/unit/} und \texttt{tests/integration/} aufgeteilt.
            Die Testfälle werden über \texttt{gtest\_discover\_tests()} automatisch bei CTest registriert.
            \sig{Adrian Gabriel Axtmann}

            \subsubsection{CI-Pipeline}\label{subsubsec:impl-medusa-ci-pipeline}
                \index{CI!GitLab}
                Die \gls{ci-pipeline} läuft auf GitLab~CI in einem Docker-Image mit vorgebauten Abhängigkeiten.

                \begin{table}[ht]
                    \centering
                    \caption[Stages der GitLab-CI-Pipeline]{%
                        Stages der Pipeline und ihre Verantwortlichkeit%
                    }
                    \label{tab:impl-medusa-ci-stages}
                    \renewcommand{\arraystretch}{1.2}
                    \begin{tabular}{@{} l p{8.0cm} @{}}
                        \toprule
                        \textbf{Stage} & \textbf{Verantwortlichkeit} \\
                        \midrule
                        \texttt{prepare}        & Baut bei Bedarf das CI-Docker-Image \\
                        \texttt{check}          & Statische Prüfungen (clang-format, clang-tidy) \\
                        \texttt{build\_and\_test} & Übersetzung und GoogleTest-Suite \\
                        \texttt{deploy}         & Doxygen-Dokumentation zu GitLab~Pages \\
                        \texttt{tag}            & Automatische Release-Erzeugung \\
                        \bottomrule
                    \end{tabular}
                \end{table}

                Der Compiler-Cache \texttt{ccache} reduziert die Übersetzungszeit bei Folge-Jobs.
                Die Slicer-Baseline besteht aus drei STL-Dateien: \texttt{Cube.stl} (prismatisch), \texttt{T.stl} (T-förmig mit Überhängen) und \texttt{Bipyramide.stl} (gekrümmte Flächen).
                \sig{Adrian Gabriel Axtmann}

        \subsection{Implementierung der Slicing-Strategien}\label{subsec:impl-medusa-slicing-strategies}
            Alle drei implementierten Strategien sind im Modul \texttt{medusa::slicing} vereint.
            Die Bibliothek bindet \texttt{medusa::core} und \texttt{medusa::geometry} als öffentliche Abhängigkeiten ein und hält Eigen und libigl als private Abhängigkeiten.
            \sig{Adrian Gabriel Axtmann}

            \subsubsection{Gemeinsame Datenstrukturen und Schnittstellen}\label{subsubsec:impl-medusa-slicing-common}
                \index{Slicing!gemeinsame Datenstrukturen}
                Das Strategie-Interface \texttt{ISlicer} definiert drei Methoden.

                \begin{lstlisting}[style=cppStyle, caption={Strategie-Interface \texttt{ISlicer}.}, label={lst:impl-medusa-islicer}]
class ISlicer
{
public:
    virtual ~ISlicer() = default;

    virtual std::vector<Layer> slice(const geometry::TriangleMesh& mesh) = 0;
    virtual std::string        name()  const = 0;
    virtual float              coverage() const { return 0.0f; }
};
                \end{lstlisting}

                Die Schichtrepräsentation \texttt{Layer} enthält Konturpunkte als flache Folge, Normalen, Face-IDs, Schichtdicke, Aufbauachse und einen Trennindex zwischen Wand- und Infill-Segmenten.

                \begin{lstlisting}[style=cppStyle, caption={Schichtrepräsentation \texttt{Layer}.}, label={lst:impl-medusa-layer}]
struct Layer
{
    std::uint32_t                index;
    std::uint32_t                branch_id;
    std::vector<glm::vec3>       contour_points;
    std::vector<glm::vec3>       contour_normals;
    std::vector<std::uint32_t>   contour_face_ids;
    float                        thickness;
    int                          up_axis;
    std::size_t                  contour_count;
};
                \end{lstlisting}

                Die Pipeline-Orchestrierung erfolgt über \texttt{SlicingPipeline::execute()}, das die Post-Slicer-Stufen (\gls{infill}-Generierung, Bewegungsplanung) verkettet.
                Für den Kristall-Slicer gilt eine Ausnahme.
                Liefert er bereits einen nicht-leeren Toolpath zurück, überspringt die Pipeline die generischen Post-Slicer-Stufen und verwendet den Kristall-Toolpath direkt, um die 6-\ac{DOF}-Pfade zu erhalten.

                \sig{Adrian Gabriel Axtmann}

            \subsubsection{Implementierung des planaren Algorithmus}\label{subsubsec:impl-planar-algorithm}
                \index{Slicing-Algorithmus!Planar!Implementierung}
                Die Klasse \texttt{PlanarSlicer} folgt der in \cref{subsubsec:concept-planar-core-idea} konzipierten fünfstufigen Pipeline.
                \sig{Adrian Gabriel Axtmann}

                \paragraph{Schichtpositionen}\label{para:impl-planar-layer-positions}
                    Die Schichtpositionen werden gemäß \cref{eq:concept-planar-layer-positions} mit Mittenversatz berechnet.
                    Die letzte Schichtposition wird auf $a_{\max} - 10^{-6}$ geklemmt, um numerische Stabilität zu gewährleisten.
                    \sig{Adrian Gabriel Axtmann}

                \paragraph{Ebene-Dreieck-Schnitt}\label{para:impl-planar-intersection}
                    Pro Schicht werden alle \glspl{triangle} iteriert.
                    Pro \gls{triangle} werden die vorzeichenbehafteten Abstände der \glspl{vertex} zur Schichtebene berechnet.
                    Der \emph{P3-Fix} vermeidet Doppelzählung: eine \gls{edge} wird übersprungen, sobald ihr Endpunkt exakt auf der Ebene liegt.
                    \sig{Adrian Gabriel Axtmann}

\begin{lstlisting}[style=cppStyle, caption={Schnittberechnung mit P3-Fix (vereinfacht).}, label={lst:impl-planar-intersection}]
for (std::size_t e = 0; e < 3 && k < 2; ++e)
{
    const std::size_t a = e;
    const std::size_t b = (e + 1) % 3;

    if ((d[a] > 0.0f) != (d[b] > 0.0f))
    {
        if (std::abs(d[b]) >= k_plane_eps)  // P3-Fix: b auf Ebene -> naechste Kante
        {
            const float t = d[a] / (d[a] - d[b]);
            pts[k++] = v[a] + t * (v[b] - v[a]);
        }
    }
    else if (std::abs(d[a]) < k_plane_eps)  // a exakt auf Ebene
    {
        pts[k++] = v[a];
    }
}
\end{lstlisting}

                \paragraph{Konturverkettung}\label{para:impl-planar-chaining}
                    Die gierige \gls{nearest-neighbor}-Verkettung wählt pro Schritt das \gls{line-segment} mit dem kleinsten quadrierten Abstand zum Cursor.
                    Die \gls{snap-tolerance} beträgt $10^{-4}\,\si{mm}$.
                    Die Reihenfolgenoptimierung startet beim Ring mit der kleinsten Z-Koordinate des Schwerpunkts.
                    \sig{Adrian Gabriel Axtmann}

            \subsubsection{Implementierung des Base-Algorithmus}\label{subsubsec:impl-base-algorithm}
                \index{Slicing-Algorithmus!Base!Implementierung}
                Die Klasse \texttt{BaseSlicer} folgt der in \cref{subsubsec:concept-base-core-idea} konzipierten queue-getriebenen Mehrlauf-Architektur.
                Eine \texttt{std::queue<SliceRun>} hält anstehende Läufe.
                Ein Iterationslimit von 20 schützt gegen Spawning-Kaskaden.
                \sig{Adrian Gabriel Axtmann}

                \paragraph{Lokales UV-Frame}\label{para:impl-base-uv-frame}
                    Die Konstruktion erfolgt über ein \gls{gram-schmidt}.
                    Falls $\vec{G}$ nahezu parallel zur Y-Achse liegt, wird die X-Achse als Referenz verwendet.

                    \begin{lstlisting}[style=cppStyle, caption={Konstruktion des lokalen UV-Frames.}, label={lst:impl-base-uvframe}]
void BaseSlicer::computeUV(const glm::vec3& growth, glm::vec3& u, glm::vec3& v)
{
    glm::vec3 ref = (std::abs(glm::dot(growth, glm::vec3(0.0f, 1.0f, 0.0f))) > 0.9f)
                    ? glm::vec3(1.0f, 0.0f, 0.0f)
                    : glm::vec3(0.0f, 1.0f, 0.0f);
    u = glm::normalize(glm::cross(growth, ref));
    v = glm::normalize(glm::cross(u, growth));
}
                    \end{lstlisting}

                \paragraph{Überhangerkennung}\label{para:impl-base-overhang}
                    Pro Schicht wird das \gls{bounding-rectangle} in der \gls{uv-plane} berechnet.
                    Überschreitet die Differenz zur Vorgängerschicht die Schwelle
                    $\text{thickness} \cdot \tan(\text{overhang\_angle})$, wird ein neuer Lauf gespawnt.
                    Markierungsflags verhindern wiederholtes Spawning.
                    \sig{Adrian Gabriel Axtmann}

                    \begin{lstlisting}[style=cppStyle, caption={Überhangerkennung und Branch-Spawning in \texttt{+U}-Richtung (analog für $-\vec{U}$, $\pm\vec{V}$)}, label={lst:impl-base-overhang}]
const float tolerance = thickness *
    std::tan(glm::radians(std::clamp(mParams.overhang_angle,
                                     BaseParams::OVERHANG_ANGLE_MIN,
                                     BaseParams::OVERHANG_ANGLE_MAX)));

if (!queuedUPlus &&
    currentBounds.uMax > prevBounds.uMax + tolerance)
{
    allowed.uMax  = prevBounds.uMax;
    pending.push({u, prevBounds.uMax, nextBranchId++});
    queuedUPlus   = true;
}
                    \end{lstlisting}

                \paragraph{Liang-Barsky-Clipping}\label{para:impl-base-clipping}\index{Liang-Barsky-Verfahren}
                    Das Beschneiden der Segmente verwendet das \gls{liang-barsky}-Verfahren, das Segmente exakt an der Bereichsgrenze zerschneidet.
                    Dies löste das Problem des schleichenden Auswanderns.
                    \sig{Adrian Gabriel Axtmann}

            \subsubsection{Implementierung des Kristall-Algorithmus}\label{subsubsec:impl-crystal-algorithm}\index{Slicing-Algorithmus!Crystal!Implementierung}
                Die Implementierung folgt der in \cref{subsubsec:concept-crystal-core-idea} beschriebenen achtstufigen Pipeline und stützt sich auf \texttt{Eigen} und \texttt{libigl}.
                \sig{Adrian Gabriel Axtmann}

                \begin{table}[ht]
                    \centering
                    \caption[Klassen der Crystal-Implementierung]{%
                        Klassen der Crystal-Implementierung%
                    }
                    \label{tab:impl-crystal-classes}
                    \renewcommand{\arraystretch}{1.2}
                    \begin{tabular}{@{} l p{7.0cm} @{}}
                        \toprule
                        \textbf{Klasse} & \textbf{Verantwortlichkeit} \\
                        \midrule
                        \texttt{CrystalSlicer}     & Pipeline-Orchestrator, implementiert \texttt{ISlicer} \\
                        \texttt{MeshRefiner}       & Mittelpunkt-Unterteilung \\
                        \texttt{ScalarField}       & Berechnung von $\Phi$ und $\nabla\Phi$ \\
                        \texttt{IsoExtractor}      & Adaptive Iso-Flächen-Extraktion \\
                        \texttt{BranchTracker}     & Topologie-Tracking \\
                        \texttt{CrystalInfill}     & Nicht-planare Infill-Generierung \\
                        \texttt{AlternatingScheduler} & Round-Robin-Schichtreihenfolge \\
                        \bottomrule
                    \end{tabular}
                \end{table}

                \paragraph{Mesh-Verfeinerung}\label{para:impl-crystal-mesh-refinement}
                    Die Verfeinerung erfolgt iterativ über \texttt{igl::upsample()} bis die maximale \gls{edge}-Länge eine Schwelle von $1{,}0\,\si{mm}$ unterschreitet.
                    Eine Sicherheitsregel erzwingt mindestens eine Unterteilung bei weniger als 200 Vertizes.
                    \sig{Adrian Gabriel Axtmann}

                \paragraph{Skalarfeld}\label{para:impl-crystal-solve}
                    Der diskrete Laplace-Operator wird über \texttt{igl::cotmatrix()} aufgebaut.
                    Die Randvertizes für die Dirichlet-Bedingungen stammen aus den unteren und oberen $2\,\%$ der Bauteilhöhe.
                    Das lineare System wird über \texttt{igl::min\_quad\_with\_fixed} mit LDLT-Faktorisierung gelöst.
                    \sig{Adrian Gabriel Axtmann}

                    \begin{lstlisting}[style=cppStyle, caption={Aufbau und Lösung des Skalarfeld-Systems}, label={lst:impl-crystal-solve}]
Eigen::SparseMatrix<double> L;
igl::cotmatrix(V, F, L);
Eigen::SparseMatrix<double> Q = -L;  // min_quad erwartet pos. semidef.

Eigen::VectorXi known(bottom_idx.size() + top_idx.size());
Eigen::VectorXd Y(known.size());
Eigen::Index k = 0;
for (int idx : bottom_idx) { known(k) = idx; Y(k) = 0.0; ++k; }
for (int idx : top_idx)    { known(k) = idx; Y(k) = 1.0; ++k; }

igl::min_quad_with_fixed_data<double> data;
igl::min_quad_with_fixed_precompute(Q, known, {}, /*pd=*/true, data);

Eigen::VectorXd phi;
igl::min_quad_with_fixed_solve(data,
                               Eigen::VectorXd::Zero(V.rows()),
                               Y, {}, phi);
                    \end{lstlisting}

                \paragraph{Adaptive Iso-Extraktion}\label{para:impl-crystal-iso}
                    Die Schrittweite wird adaptiv an den mittleren Gradienten angepasst,
                    geklemmt zwischen $0{,}05$ und $8{,}0$ der uniformen Referenzschrittweite.
                    Ein Iterationslimit von 2000 Iso-Schritten dient als Schutz.
                    \sig{Adrian Gabriel Axtmann}

                    \begin{lstlisting}[style=cppStyle, caption={Adaptive Schrittweite der Iso-Extraktion.}, label={lst:impl-crystal-adaptive}]
// Einmalig vor der Schleife: Referenzschritt und Klemmbounds
const float dphi_uniform = layer_thickness / h_range
                         * (field.phi_max - field.phi_min);
const float dphi_min = dphi_uniform * cfg.dphi_min_factor;  // Faktor 0,05
const float dphi_max = dphi_uniform * cfg.dphi_max_factor;  // Faktor 8,0

// Am Ende jeder Iso-Iteration: lokalen Gradienten messen und klemmen
const float g_use     = (g_count > 0)
                        ? float(g_sum / g_count) : g_seed;
const float dphi_next = std::clamp(layer_thickness * g_use,
                                   dphi_min, dphi_max);
phi_cur += dphi_next;
                    \end{lstlisting}

                \paragraph{Branch-Tracking}\label{para:impl-crystal-branch-tracking}
                    \index{Zweig!Tracking}
                    Iso-Komponenten werden über Schwerpunktabstände zugeordnet.
                    Die Drift-Schwelle beträgt $\delta_\mathrm{drift} = \max(0{,}5\,\si{mm},\,4t)$.
                    Split- und Merge-Ereignisse werden erkannt und protokolliert.
                    \sig{Adrian Gabriel Axtmann}

                \paragraph{Infill-Generierung}\label{para:impl-crystal-infill}
                    Die Tangentialbasis wird über \ac{PCA} bestimmt.
                    Bei isotropen Punktwolken ($\lambda_2/\lambda_1 > 0{,}7$) wird die globale X-Achse als Fallback verwendet.
                    Der Scanline-Sweep erzeugt ein \gls{snake-pattern}.
                    Die \gls{idw}-Anhebung mit $k = 6$ Nachbarn projiziert die Punkte auf die Iso-Fläche.
                    \sig{Adrian Gabriel Axtmann}

                \paragraph{Werkzeugausrichtung}\label{para:impl-crystal-pose}
                    Die implementierte Ausrichtung beschränkt sich auf die Werkzeug-Z-Achse $\vec{z}_\mathrm{tool}$.
                    Pro Segment wird das arithmetische Mittel der Oberflächennormalen an beiden Endpunkten
                    normiert und als \texttt{orientation}-Feld der \texttt{Segment}-Struktur gespeichert.

                    \begin{lstlisting}[style=cppStyle, caption={Berechnung der Werkzeugausrichtung}, label={lst:impl-crystal-tool-frame}]
glm::vec3 n = nrm[i] + nrm[i + 1];
const float l = glm::length(n);
s.orientation = (l > 1e-8f) ? (n / l) : glm::vec3(0, 1, 0);
                    \end{lstlisting}

                    Die vollständige Orientierung als \gls{quaternion} (Felder \texttt{qw}, \texttt{qx}, \texttt{qy},
                    \texttt{qz} der \texttt{Waypoint}-Struktur) ist eine offene Aufgabe; das aktuelle System
                    überträgt keinen vollständigen Drehrahmen an \ac{kronos}.
                    \sig{Adrian Gabriel Axtmann}

            \subsubsection{Performance-Charakterisierung}\label{subsubsec:impl-slicing-performance}
                \index{Slicing!Performance}
                Beim planaren Algorithmus dominiert die Schnittstufe mit $\mathcal{O}(NT)$.
                Beim Base"=Algorithmus multipliziert sich dies mit der Branch"=Anzahl zu $\mathcal{O}(BNT)$.
                Beim Crystal-Algorithmus dominiert die Solver-Stufe mit etwa $\mathcal{O}(V^{1{,}5})$.
                Für die Validierungsgeometrien liegen alle Laufzeiten im einstelligen Sekundenbereich.
                \sig{Adrian Gabriel Axtmann}

        \subsection{Zusammenfassung}\label{subsec:impl-medusa-summary}
            Die Implementierung setzt alle Schichten der konzeptionellen Architektur um.
            Die durchgängige Verarbeitungskette vom Mesh-Import über Geometrieaufbereitung, Slicing, \gls{infill}-Generierung und Bewegungsplanung bis zum atomaren \gls{json}-Export funktioniert in allen drei Strategien.
            Die Visualisierung mit Mesh-, Hilfsgeometrie-, Toolpath- und Feld-Overlay-Schicht sowie die Schnittstelle zu \ac{kronos} sind lauffähig.
            \sig{Adrian Gabriel Axtmann}

            \paragraph{Bekannte Einschränkungen}\label{para:impl-medusa-summary-limitations}
                Vier Klassen von Einschränkungen bestehen.
                Im Kristall"=Algorithmus fehlen Mehrwand-Konturen, krümmungsadaptive Schichtdicke und Branch-Kollisionserkennung.
                Der Sphären"=Algorithmus ist nicht implementiert.
                Die Eingangsstufe führt keine Mesh-Reparatur durch.
                Die plattformübergreifende Validierung ist nur unter Linux automatisiert.
                \sig{Adrian Gabriel Axtmann}

            \paragraph{Übergang zur Evaluation}\label{para:impl-medusa-summary-transition}
                Die dokumentierte Implementierung ist die Grundlage für die in \cref{chap:evaluation} folgende Bewertung.
                Dort wird die Funktionsfähigkeit der drei Slicing-Strategien geprüft, die Ausführbarkeit der \glspl{toolpath} durch \ac{kronos} bewertet und das Gesamtsystem gegen die Forschungsfragen reflektiert.
                \sig{Adrian Gabriel Axtmann}

    \section{Implementierung der Komponente KRONOS}\label{sec:implementation-of-component-kronos}
        \subsection{Systemarchitektur}\label{subsec:kronos-system-architecture}
            \subsubsection{Virtualisierung mittels Docker}
            Die Implementierung von \ac{kronos} erfolgt in einer containerisierten Entwicklungsumgebung.
            Das Docker-Image baut auf \texttt{osrf/ros:humble-desktop-full} auf und installiert darauf die für den Betrieb benötigten ROS-2-, MoveIt- und Universal-Robots-Pakete.
            Zusätzlich werden Build-Werkzeuge wie \texttt{colcon}, \texttt{cmake} und \texttt{python3-rosdep} sowie die für inverse Kinematik benötigten Bibliotheken installiert.
            TRAC-IK wird dabei nicht ausschließlich als Binärpaket eingebunden, sondern im Container aus dem Quelltext gebaut und an ROS~2 Humble angepasst.
            Dazu war ein quellcode-seitiger Patch notwendig: Der im Rolling-Branch von TRAC-IK verwendete Header \texttt{urdf/model.hpp} existiert unter Humble nicht mehr und musste durch \texttt{urdf/model.h} ersetzt werden.
            Diese Kapselung reduziert den Einrichtungsaufwand auf dem Host-System deutlich und stellt sicher,
            dass Entwicklung, Build und Ausführung mit identischen Abhängigkeiten erfolgen.
            \sig{Nik Wachsmann}

            \subsubsection{Subkomponenten}
            Innerhalb des ROS-2-Workspaces ist \ac{kronos} als Paket \texttt{ur3e\_control} umgesetzt.
            Die Implementierung gliedert sich in zwei operative Teilbereiche sowie einen ergänzenden Versuchsknoten.

            \paragraph{Simulationsumgebung}
            Die Simulationsumgebung wird über die Launch-Datei \texttt{ur3e\_full.launch.py} gestartet.
            Zunächst wird der \texttt{ur\_robot\_driver} im Modus \textit{fake hardware} gestartet,
            sodass die Steuerungsschnittstellen des UR3e ohne reale Hardware verfügbar sind.
            Anschließend werden zeitverzögert die MoveIt-Konfiguration und RViz gestartet,
            damit der Treiber seine Controller vollständig initialisieren kann.
            Die Visualisierung erfolgt über eine vorbereitete RViz-Konfiguration,
            welche den Roboterzustand, die Planungsszene und später auch die vom Druckprozess erzeugten Marker darstellt.
            
            \paragraph{Trajektoriensteuerung}
            Die eigentliche Ausführungslogik befindet sich im Knoten \texttt{ur3e\_printer}.
            Dieser liest die von Medusa erzeugten Pfaddaten ein,
            transformiert sie bei Bedarf aus dem Werkstück-Koordinatensystem in das Roboterkoordinatensystem
            und erzeugt daraus MoveIt-kompatible Zielposen.
            Als Kinematik-Plugin ist dabei \texttt{kdl\_kinematics\_plugin/KDLKinematicsPlugin} konfiguriert;
            TRAC-IK steht zwar im Container zur Verfügung, wird jedoch nur im ergänzenden Versuchsknoten genutzt.

            \paragraph{Experimentalknoten}
            Der Knoten \texttt{ur3e\_trac\_ik\_controller} löst inverse Kinematik direkt über die TRAC-IK-Bibliothek
            und führt die resultierenden Gelenkwinkel über MoveIt aus.
            Er dient ausschließlich der isolierten Erprobung des TRAC-IK-Solvers
            und ist nicht in den regulären Druckablauf eingebunden.
            \sig{Nik Wachsmann}

            \subsubsection{Komponenten Kommunikation}
            Die Kommunikation zwischen den Komponenten erfolgt vollständig innerhalb des ROS~2 Ökosystems.
            Der Druckknoten erhält die \url{robot\_description} und die Kinematikparameter beim Start aus der Launch Datei
            und erzeugt daraus ein \texttt{Move\-Group\-Interface} für die Planungsgruppe \url{ur\_manipulator}.
            Bewegungspläne werden nicht manuell an Controller serialisiert,
            sondern über die MoveIt-Planungsschnittstelle an den aktiven \url{joint\_trajectory\_controller} übergeben.
            Ergänzend publiziert der Knoten ein persistentes Marker-Array auf dem Topic \url{printed\_part},
            um bereits ausgeführte Drucksegmente in RViz sichtbar zu machen.
            \sig{Nik Wachsmann}

        \subsection{Entwicklungs- und Laufzeitumgebung}
            Die Entwicklungsumgebung ist so aufgebaut,
            dass der Quelltext des Pakets \url{ur3e\_control} vom Host-System in den Container eingebunden wird,
            während Build-, Installations- und Log-Verzeichnisse als Docker-Volumes persistent gespeichert werden.
            Dadurch können Änderungen lokal mit einem Editor vorgenommen und im Container unmittelbar mit colcon neu gebaut werden.
            Für die grafischen Werkzeuge wird der X11-Socket des Hosts in den Container durchgereicht.
            Ein Hilfsskript bereitet die notwendige Freigabe für Docker oder Podman vor.

            Eine getrennte Laufzeitumgebung wurde nicht implementiert.
            Stattdessen dient dieselbe Containerkonfiguration sowohl für die Entwicklung als auch für die Ausführung in der Simulation.
            Dieses Vorgehen ist für einen Prototypen zweckmäßig,
            weil Integrationsprobleme zwischen Build- und Zielumgebung vermieden werden.
            Über die Launch-Dateien lässt sich dabei zwischen verschiedenen Betriebsarten unterscheiden:
            Die Standardkonfiguration startet den simulierten UR3e mit MoveIt und RViz,
            während der Anschluss an reale Hardware über den offiziellen \texttt{ur\_robot\_driver} prinzipiell vorbereitet ist.
            \sig{Nik Wachsmann}

        \subsection{Start- und Konfigurationslogik}
            Die Initialisierung von \ac{kronos} erfolgt bewusst nicht direkt im C++-Knoten,
            sondern über Python-basierte ROS-2-Launch-Dateien.
            Dadurch lassen sich Robotermodell,
            semantische Beschreibung und Kinematik-Konfiguration bereits vor dem Start des Druckknotens zusammenführen.
            Für den eigentlichen Druckprozess ist insbesondere \texttt{ur3e\_printer.launch.py} relevant,
            da dort die für MoveIt notwendigen Parameter sowie der Pfad zur Eingabedatei an den Knoten \texttt{ur3e\_printer} übergeben werden.
            Über eine paketinterne \texttt{kinematics.yaml} wird das Standard-Kinematik-Plugin durch
            \texttt{kdl\_kinematics\_plugin/KDLKinematicsPlugin} ersetzt;
            konfiguriert sind eine Suchauflösung von 0{,}005\,rad, ein Solver-Timeout von 50\,ms und bis zu drei Lösungsversuche.

            Listing~\ref{lst:kronos-launch-printer-node} zeigt den zentralen Teil dieser Initialisierung.
            Sichtbar ist,
            dass der Knoten nicht nur die \gls{json}-Datei erhält,
            sondern gleichzeitig auch die vom Launch-System erzeugten Modellparameter eingebunden bekommt.

            \begin{lstlisting}[language=Python,caption={Initialisierung des Druckknotens in der Launch-Datei},label={lst:kronos-launch-printer-node}]
toolpath_file = LaunchConfiguration("toolpath_file")

printer_node = Node(
    package="ur3e_control",
    executable="ur3e_printer",
    name="ur3e_printer",
    output="screen",
    parameters=[
        robot_description,
        robot_description_semantic,
        robot_description_kinematics,
        {"toolpath_file": toolpath_file},
    ],
)
            \end{lstlisting}

            Ergänzt wird dieser Startpfad durch \texttt{ur3e\_full.launch.py},
            das zunächst den UR-Treiber im Simulationsmodus und anschließend mit einer Verzögerung von fünf Sekunden MoveIt und RViz startet.
            Diese zeitliche Staffelung ist relevant,
            weil die Controller des Treibers vor der Planung vollständig verfügbar sein müssen.
            \sig{Nik Wachsmann}

        \subsection{Einlesen und Aufbereitung der Trajektoriendaten}
            Für den beschriebenen Fertigungsablauf ist ausschließlich das von Medusa erzeugte \gls{json}-Format relevant.
            Der Knoten \texttt{ur3e\_printer} erwartet darin ein Top-Level-Feld \texttt{waypoints},
            dessen Einträge kartesische Positionen in Millimetern sowie \glspl{quaternion} zur Werkzeugorientierung enthalten.
            Beim Einlesen werden diese Werte zunächst validiert und anschließend in das von ROS und MoveIt erwartete Einheitensystem in Metern überführt.
            Das optionale Top-Level-Feld \texttt{reference\_frame} steuert dabei, ob eine Koordinatentransformation stattfindet:
            Der Wert \texttt{workpiece} aktiviert die Umrechnung vom Werkstück- ins Roboterbasissystem;
            fehlt das Feld oder enthält es \texttt{base\_link}, werden die Werte unverändert übernommen.

            Ein zusätzlicher Verarbeitungsschritt ist notwendig,
            wenn Medusa die Bahn nicht im globalen Roboterkoordinatensystem,
            sondern relativ zum Werkstück exportiert.
            In diesem Fall transformiert \ac{kronos} die Positionen von der Y-up-Darstellung des Werkstücks in das in ROS übliche Z-up-System.
            Parallel dazu wird aus der übergebenen Werkzeugausrichtung eine zur Roboterbasis konsistente Orientierung berechnet.
            Listing~\ref{lst:kronos-json-transform} zeigt den Kern dieser Umrechnung.

            \begin{lstlisting}[language=C++,caption={Einlesen und Transformation werkstückbezogener \gls{json}-Wegpunkte},label={lst:kronos-json-transform}]
double x = waypoint.at("x").get<double>() / 1000.0;
double y = waypoint.at("y").get<double>() / 1000.0;
double z = waypoint.at("z").get<double>() / 1000.0;

if (is_workpiece_frame) {
  const tf2::Vector3 robot_position = medusa_workpiece_position_to_robot_basis(x, y, z);
  x = robot_position.x();
  y = robot_position.y();
  z = robot_position.z();

  const tf2::Vector3 medusa_tcp_direction = tf2::quatRotate(medusa_q, tf2::Vector3(0.0, 0.0, 1.0));
  const tf2::Vector3 robot_tcp_direction = medusa_tcp_direction_to_robot_basis(medusa_tcp_direction);
  robot_q = quaternion_from_z_axis(robot_tcp_direction);
}
            \end{lstlisting}

            Die aufbereiteten Wegpunkte werden intern nicht als \glspl{quaternion} weiterverarbeitet,
            sondern in Roll-, Pitch- und Yaw-Winkel überführt und in einer flachen Datenstruktur gespeichert.
            Diese Normalisierung reduziert die Komplexität der nachfolgenden Pose-Erzeugung,
            da jede geplante Zielpose aus einer einheitlichen Repräsentation zusammengesetzt werden kann.
            \sig{Nik Wachsmann}

        \subsection{Koordinatentransformation und Orientierungslogik}
            \ac{kronos} führt das werkstückbezogene Koordinatensystem der Eingabedatei,
            das Basissystem des Roboters und die Pose des Tool Center Points in eine gemeinsame Repräsentation über.
            Dies ist notwendig,
            weil Medusa die Druckbahn relativ zum Bauteil beschreibt,
            während MoveIt Zielposen in der Roboterbasis erwartet.
            Besonders relevant ist dabei die Umrechnung von der Medusa-Y-up-Darstellung in das in ROS übliche Z-up-System,
            da hiervon sowohl die Positionen als auch die Bedeutung der Werkzeugorientierung abhängen.

            Die in der \gls{json}-Datei enthaltenen \glspl{quaternion} werden deshalb nicht direkt übernommen,
            sondern zunächst zur Bestimmung der lokalen TCP-Richtung ausgewertet.
            Diese Richtung wird in das Roboterbasissystem transformiert und anschließend wieder in eine gültige Zielorientierung überführt.
            Ergänzend werden Mehrdeutigkeiten in der Rotationsdarstellung reduziert,
            indem Orientierungen intern vereinheitlicht und Gelenkwinkel vor der Ausführung entrollt werden.
            Dadurch bleiben Werkzeughaltung und Trajektorienverlauf entlang benachbarter Wegpunkte konsistent,
            was sowohl die Erreichbarkeit als auch die visuelle Plausibilitätsprüfung in RViz verbessert.
            \sig{Nik Wachsmann}

        \subsection{Bahnplanung und Ausführung}
            Nach dem erfolgreichen Laden der \gls{json}-Datei initialisiert der Knoten die MoveIt-Schnittstelle,
            wartet auf gültige Gelenkzustände und berechnet anschließend zunächst eine sichere Annäherung an den ersten Wegpunkt.
            Dazu wird oberhalb des ersten Zielpunkts eine Annäherungspose erzeugt,
            für die mehrere Planungsversuche durchgeführt werden.
            Von diesen Kandidaten wird der kürzeste Gelenkpfad ausgewählt und ausgeführt.
            Der letzte Abschnitt bis zum tatsächlichen Startpunkt wird bevorzugt als kartesische Gerade geplant,
            um den Eintritt in den Druckpfad kontrolliert und reproduzierbar zu halten.

            Der eigentliche Druckpfad wird blockweise mit maximal 50 Wegpunkten verarbeitet.
            Für jeden Block ruft \ac{kronos} \texttt{computeCartesianPath} auf und bewertet anschließend,
            welcher Anteil der angeforderten Punkte tatsächlich planbar war.
            Der Sprungfilter (\texttt{jump\_threshold}) ist auf 0{,}0 gesetzt und damit deaktiviert,
            da er im Simulationsbetrieb mit Fake-Hardware zu Fehlalarmen führt.
            Wird die Mindestabdeckung unterschritten,
            versucht das System lokal eine Wiederaufnahme im Gelenkraum,
            statt den gesamten Auftrag sofort zu verwerfen.
            Listing~\ref{lst:kronos-batch-planning} zeigt den Kern dieser Steuerungslogik.

            \begin{lstlisting}[language=C++,caption={Batchweise kartesische Planung und Coverage-Auswertung},label={lst:kronos-batch-planning}]
moveit_msgs::msg::RobotTrajectory trajectory;
double coverage = move_group.computeCartesianPath(poses, CART_EEF_STEP, CART_JUMP_THRESH, trajectory);

size_t planned_pts = static_cast<size_t>(std::round(coverage * static_cast<double>(poses.size())));

if (coverage < MIN_COVERAGE) {
  if (planned_pts == 0) {
    // lokale Recovery-Planung im Gelenkraum
  }
}
            \end{lstlisting}

            Vor der Ausführung einer Trajektorie werden zwei zusätzliche Verarbeitungsschritte vorgenommen.
            Erstens ergänzt \ac{kronos} mit der \textit{Time-Optimal Trajectory Generation} fehlende Zeitstempel,
            Geschwindigkeiten und Beschleunigungen,
            da \texttt{compute\-Cartesian\-Path} zunächst nur diskrete Gelenkpositionen liefert.
            Zweitens werden eventuelle $2\pi$-Sprünge in den Gelenkwinkeln entfernt,
            damit äquivalente Winkelrepräsentationen nicht zu unnötigen Umorientierungen des Roboters führen.
            Erst danach wird die Trajektorie über MoveIt an den \texttt{joint\_trajectory\_controller} übergeben.
            \sig{Nik Wachsmann}

        \subsection{Recovery-Strategien und Fehlerbehandlung}
            \ac{kronos} behandelt unvollständige kartesische Planung nicht sofort als Abbruch,
            sondern bewertet nach jedem Aufruf von \texttt{computeCartesianPath},
            welcher Anteil eines Wegpunktblocks tatsächlich planbar war.
            Wird die Mindestabdeckung unterschritten,
            können verwertbare Teiltrajektorien weiterhin ausgeführt werden;
            bei vollständig unplanbaren Abschnitten versucht das System stattdessen eine lokale Wiederaufnahme im Gelenkraum.
            Dazu werden bis zu zehn nachfolgende Wegpunkte als Wiederaufnahmekandidaten geprüft;
            der erste erfolgreich planbare Punkt wird als Einstieg gewählt und übersprungene Wegpunkte werden protokolliert.
            Dadurch lassen sich einzelne problematische Segmente überbrücken,
            ohne den gesamten Fertigungsauftrag unmittelbar zu verwerfen.

            Harte Fehlerzustände führen dennoch zu einem sofortigen Stopp.
            Dazu zählen insbesondere ungültige Eingabedaten,
            fehlgeschlagene Initialisierungsschritte,
            nicht verfügbare Gelenkzustände oder eine nicht planbare Annäherung an den ersten Wegpunkt.
            Ergänzend protokolliert \ac{kronos} solche Situationen zur Laufzeit und visualisiert den bis dahin erfolgreich ausgeführten Pfad in RViz.
            Damit wird zwischen tolerierbaren Teilfehlern und Fehlern unterschieden,
            bei denen keine belastbare oder fachlich sinnvolle Ausführung mehr gewährleistet ist.
            \sig{Nik Wachsmann}

        \subsection{Visualisierung und Laufzeitrückmeldung}
            Neben der reinen Bewegungsausführung implementiert \ac{kronos} auch eine explizite Rückmeldung des Druckfortschritts.
            Zu diesem Zweck sammelt der Knoten bereits erfolgreich abgefahrene Segmente und publiziert sie als \texttt{MarkerArray} auf dem Topic \texttt{printed\_part}.
            In RViz entsteht dadurch schrittweise eine visuelle Repräsentation des bislang ausgeführten Druckpfads.

            Diese Visualisierung erfüllt zwei Funktionen.
            Zum einen erlaubt sie eine direkte Plausibilitätsprüfung während der Entwicklung,
            etwa wenn kartesische Teilpfade unerwartet abbrechen oder übersprungen werden.
            Zum anderen unterstützt sie die Fehlersuche,
            weil sich geplante und bereits ausgeführte Bahnsegmente klar voneinander abgrenzen lassen.
            Die Visualisierung ist damit kein bloßes Zusatzmerkmal,
            sondern ein integraler Bestandteil der prototypischen Inbetriebnahme.
            \sig{Nik Wachsmann}
        
        \subsection{Einbindung zentraler Frameworks}
            \ac{kronos} verwendet mehrere Frameworks,
            die jeweils einen klar abgegrenzten Teil der Steuerungs- und Simulationslogik übernehmen.
            \paragraph{ROS~2}
            ROS~2 bildet die technische Grundlage der gesamten Implementierung.
            Das Paket \texttt{ur3e\_control} ist als reguläres \texttt{ament\_cmake}-Paket umgesetzt und bindet darüber unter anderem \texttt{rclcpp},
            Nachrichten- und Geometriebibliotheken sowie die Schnittstellen von MoveIt ein.
            ROS~2 übernimmt damit sowohl das Starten der Knoten über Launch-Dateien als auch die Verteilung von Parametern,
            Zustandsinformationen und Visualisierungsdaten.

            \paragraph{RViz}
            RViz wird als primäres Werkzeug zur Laufzeitbeobachtung eingesetzt.
            Innerhalb der Simulation visualisiert es das Roboter-Modell,
            die durch MoveIt verwaltete Planungsszene und den Fortschritt des Druckpfads.
            Die bereits abgefahrenen Segmente werden vom Druckknoten als \texttt{LINE\_STRIP}-Marker publiziert,
            wodurch sich die tatsächliche Ausführung schrittweise nachvollziehen lässt.

            \paragraph{MoveIt}
            MoveIt übernimmt die eigentliche Bewegungsplanung.
            Als Kinematik-Plugin wird \texttt{kdl\_kinematics\_plugin/KDLKinematicsPlugin} verwendet.
            Die Library \texttt{TRAC-IK} wird zwar im Container bereitgestellt, ist aber ausschließlich im Versuchsknoten \url{ur3e_trac_ik_controller} aktiv.
            Im Knoten \texttt{ur3e\_printer} wird zunächst eine Annäherungsbewegung an den ersten Wegpunkt geplant,
            wobei zehn Planungsversuche erzeugt und anschließend der kürzeste Gelenkpfad ausgewählt wird.
            Der Übergang zum ersten Druckpunkt erfolgt bevorzugt als kartesische Gerade;
            schlägt dies fehl, wird auf eine gelenkraumorientierte Ausweichplanung zurückgegriffen.
            Der restliche Druckpfad wird in Blöcke von jeweils bis zu 50 Wegpunkten aufgeteilt und über \texttt{computeCartesianPath} geplant.
            Für jede Teiltrajektorie wird eine Mindestabdeckung von 95\,\% gefordert.
            Unvollständige Abschnitte werden entweder teilweise ausgeführt oder,
            falls keine kartesische Lösung vorliegt,
            durch eine lokale Recovery-Planung im Gelenkraum überbrückt.
            Anschließend wird auf die geplanten Gelenkpositionen eine zeitoptimale Parametrisierung angewendet,
            damit der \texttt{joint\_trajectory\_controller} die Trajektorie mit gültigen Zeitstempeln,
            Geschwindigkeiten und Beschleunigungen ausführen kann.

            Für die Orientierung des Werkzeugs unterstützt der Druckknoten direkt übergebene \glspl{quaternion} im \ac{kronos}-\gls{json}-Format.
            Zusätzlich kann ein werkstückbezogenes Koordinatensystem verwendet werden,
            das zur Laufzeit in das ROS-übliche Z-up-System transformiert wird.
            Zur Vermeidung von Gelenksprüngen werden geplante Trajektorien vor der Ausführung entrollt,
            indem $2\pi$-Umschläge relativ zum aktuellen Gelenkzustand entfernt werden.
            Damit adressiert die Implementierung direkt ein typisches Problem kontinuierlicher Roboterbahnen,
            bei dem äquivalente Gelenkwinkelrepräsentationen sonst zu abrupten Umorientierungen führen würden.
            \sig{Nik Wachsmann}